\documentclass[11pt]{article}

\usepackage[margin=1in]{geometry}
\usepackage{amsmath}
\usepackage{amssymb}
\usepackage{booktabs}
\usepackage[hidelinks]{hyperref}
\usepackage{url}

\title{Can 3D Gaussian Splatting Help Monte Carlo Denoising?}
\author{vkGSplat research note}
\date{Started May 4, 2026}

\begin{document}
\maketitle

\section{Executive Thesis}

As of May 4, 2026, the most defensible answer is:

\begin{quote}
3D Gaussian Splatting is probably not a strong replacement for a conventional
single-frame Monte Carlo image denoiser, but it is a credible and timely
representation for \emph{world-space or path-space radiance caching},
\emph{sample reconstruction}, and \emph{variance reduction}. The closest
published evidence is GSCache, which adapts 3D Gaussian splatting into an
online multi-level path-space radiance cache for volume path tracing.
\end{quote}

For vkGSplat, this means the research direction is worth pursuing, but the
right formulation is not ``render a noisy image, train 3DGS, render a clean
image.'' The stronger formulation is:
\[
  \begin{aligned}
  &\text{low-spp path samples}
  \rightarrow
  \text{world/path-space Gaussian cache}
  \\
  &\rightarrow
  \text{lower-variance indirect radiance estimate}
  \rightarrow
  \text{existing temporal/SVGF-style reconstruction}.
  \end{aligned}
\]

My current subjective likelihoods:

\begin{itemize}
  \item Useful internal vkGSplat experiment: \textbf{high}, about 0.75.
  \item Publishable prototype for static or slowly changing scenes: \textbf{medium-high}, about 0.6.
  \item General drop-in denoiser for arbitrary path-traced frames: \textbf{medium-low}, about 0.35.
  \item Unbiased estimator with clean theory and real-time performance: \textbf{low-medium}, about 0.25 unless framed as a control variate with careful correction.
\end{itemize}

\section{Why This Is Subtle}

Monte Carlo denoising is usually an image/sample reconstruction problem:
given noisy estimates of pixel radiance plus auxiliary buffers such as depth,
normal, albedo, motion, variance, primitive ID, and per-sample features, recover
a low-error image. 3DGS is instead a scene/radiance-field representation:
it stores a set of anisotropic kernels in 3D, projects them, and blends them.
That mismatch is the main risk.

The opportunity is that Monte Carlo path samples are naturally not just pixels.
They contain hit positions, directions, throughput, PDFs, BSDF information,
visibility events, primitive identity, and path depth. If we preserve those
sample records, a Gaussian representation can accumulate radiance in the space
where the signal is smoother than it is in screen space.

\section{Relevant Research Landscape}

\subsection{Classical and Neural Monte Carlo Denoising}

SVGF is the relevant real-time baseline: it temporally accumulates one-path-per
pixel input and drives a hierarchical image-space filter with luminance variance
\cite{schied2017svgf}. vkGSplat already has a CPU SVGF-style baseline in
\path{include/vkgsplat/denoise.h} and \path{src/core/denoise.cpp}, which makes
this the natural control.

KPCN showed that learned kernel prediction can denoise production Monte Carlo
frames using feature buffers \cite{bako2017kpcn}. Kernel-Splatting Networks
went one level deeper and learned directly from individual Monte Carlo samples,
not only per-pixel summaries \cite{gharbi2019ksn}. This matters because a 3DGS
cache should probably be fed by individual path samples, not by already averaged
pixels.

Neural Control Variates and Neural Radiance Caching are close conceptual
ancestors. Neural Control Variates use learned approximations for variance
reduction in Monte Carlo integration \cite{muller2020ncv}. Neural Radiance
Caching trains a cache online while rendering and uses it for path-traced
global illumination with little overhead \cite{muller2021nrc}. A Gaussian cache
can be read as a more explicit, splattable, geometry-aware cousin of NRC.

\subsection{3DGS Meets Monte Carlo}

The original 3DGS paper introduced explicit anisotropic Gaussian primitives
with fast visibility-aware rendering for real-time radiance-field view synthesis
\cite{kerbl2023gs}. That work is not Monte Carlo denoising, but it established
the representation and rendering performance that make this idea attractive.

3D Gaussian Splatting as Markov Chain Monte Carlo reinterprets the set of
Gaussians as MCMC samples from an underlying scene distribution and relates
updates to stochastic gradient Langevin dynamics \cite{kheradmand2024mcmc}.
This is relevant philosophically: it says Gaussian placement/densification can
be probabilistic rather than just heuristic. It does not solve MC rendering
denoising directly.

StochasticSplats introduces an unbiased Monte Carlo estimator for the volume
rendering equation in 3DGS, avoiding depth sorting and allowing a sample count
quality/cost tradeoff \cite{kheradmand2025stochasticsplats}. This is important
because it puts MC estimators inside the 3DGS renderer itself. Again, it is not
a path-tracing denoiser, but it weakens the mental boundary between splatting
and sampling.

Unified Gaussian Primitives for Scene Representation and Rendering builds
Gaussian primitives suitable for Monte Carlo path tracing under arbitrary
illumination \cite{zhou2024unified}. It gives Gaussians transmittance and phase
function semantics, then path traces them. This is close to the representation
side of our problem: if Gaussians can be light-transport primitives, they can
also plausibly be light-transport caches.

\subsection{Closest Hit: GSCache}

GSCache is the clearest evidence that this direction is real. Bauer, Wu,
Gadirov, and Ma adapt 3D Gaussian splatting into a multi-level path-space
radiance cache for volume path tracing \cite{bauer2026gscache}. The University
of Groningen portal describes the cache as trainable on the fly and dynamically
adaptable to lighting and transfer-function changes. The abstract specifically
targets slow rendering and high pixel variance caused by Monte Carlo
integration, and reports less noisy, higher-quality images without increasing
rendering cost.

Important caveat: GSCache is for path-traced volume rendering in scientific
visualization. vkGSplat currently cares about low-sample ray-tracing seeds,
surface scenes, robotics synthetic data, and compute-first Vulkan lowering.
The result is directly encouraging but not directly transferable.

The deeper design lesson is more specific than ``use Gaussians as a cache.''
GSCache organizes the cache as multiple path-length levels: each level stores
attenuated radiance for paths of a given length or subspace of path space. It
rasterizes cache images before the path tracing pass, then updates the Gaussian
cache from noisy samples produced by the path tracer. The attractive integration
claim is that the renderer need only expose attenuated radiance and path length,
rather than a full neural-training feature stack. For vkGSplat, this suggests a
surface-scene variant with explicit levels for direct hit, first indirect bounce,
second indirect bounce, and residual long-tail illumination.

\subsection{3DGS Denoising That Is Not MC Denoising}

Recent papers use 3DGS for denoising captured multi-view or burst images.
DenoiseGS uses a Gaussian reconstruction model for burst image denoising
\cite{cheng2025denoisegs}. DenoiseSplat studies noisy multi-view scene
reconstruction with feed-forward Gaussian splatting \cite{jiang2026denoisesplat}.
Denoise-GS combines self-supervised image denoising with sparse-view 3DGS
optimization \cite{xu2026denoisegs}.

These are useful adjacent signals: they show Gaussians can impose multi-view
consistency and suppress sensor/image noise. But their noise model is not the
same as Monte Carlo path-tracing variance. We should cite them as adjacent,
not as proof.

\section{Core Feasibility Argument}

\subsection{What 3DGS Gives Us}

3DGS provides:

\begin{itemize}
  \item local anisotropic kernels in world space,
  \item fast projection/splatting into image space,
  \item differentiable or incremental optimization paths,
  \item natural support for multi-view consistency,
  \item a compact explicit cache that is more inspectable than an MLP,
  \item a structure that maps well onto CUDA/tensor-rich compute.
\end{itemize}

These are exactly the properties wanted by a renderer that accumulates
many low-spp seed frames across nearby camera poses. For robotics SDG, the
scene distribution is controlled, camera motion can be logged exactly, and
perfect G-buffers are available. This makes a world-space cache much less scary
than it would be in a game engine with arbitrary dynamic content.

\subsection{Where It Can Fail}

The hardest parts are visibility, high-frequency transport, and bias.
Indirect radiance is not only a smooth function of position. It can change
sharply with normal, outgoing direction, roughness, occlusion, light visibility,
medium state, and path depth. A naive spatial Gaussian cache will leak light
across walls, smear glossy caustics, lag behind moving lights, and double-count
transport unless the estimator is carefully designed.

\section{Promising Formulations}

\subsection{Formulation A: World-Space Gaussian Radiance Cache}

This is the main bet.

Borrow the level idea from GSCache, but reinterpret it for surfaces:

\begin{itemize}
  \item level 0: primary-hit surface attributes and directly visible radiance,
  \item level 1: first-bounce indirect diffuse or rough radiance,
  \item level 2: second-bounce and longer low-frequency residual,
  \item level R: optional residual/control-variate correction samples.
\end{itemize}

This prevents the cache from averaging path contributions with very different
variance and spatial smoothness. It also gives vkGSplat a knob: start with
level 1 only, then add deeper levels if the first experiment works.

For each path sample, record:
\[
  s_i = (x_i, n_i, \omega_o, \omega_i, L_i, \beta_i, p_i,
        \text{depth}_i, \text{primitive}_i, \sigma_i^2)
\]
where $x$ is hit position, $n$ is normal, $\omega_o$ is outgoing/view direction,
$\omega_i$ is sampled incident direction, $L_i$ is radiance or contribution,
$\beta_i$ is throughput, $p_i$ is the sampling PDF, and $\sigma_i^2$ is a local
variance estimate.

Maintain Gaussian cache entries:
\[
  G_j =
  (\mu_j, \Sigma_j, \bar{n}_j, a_j, c_j(\omega), v_j, m_j,
   \text{primitive tags})
\]
where $\mu_j$ and $\Sigma_j$ define the spatial kernel, $a_j$ is opacity or
confidence, $c_j(\omega)$ stores directional radiance, $v_j$ stores variance,
and $m_j$ stores sample count or effective mass.

At a shading point, query nearby Gaussians with a gated kernel:
\[
  w_j(x,n,\omega) =
  \exp\left(-\frac{1}{2}(x-\mu_j)^T\Sigma_j^{-1}(x-\mu_j)\right)
  \cdot g_n(n,\bar{n}_j)
  \cdot g_\omega(\omega)
  \cdot g_{\text{vis}}.
\]
The cache estimate is:
\[
  \hat{L}_{cache}(x,\omega) =
  \frac{\sum_j w_j c_j(\omega)}
       {\sum_j w_j + \epsilon}.
\]

Use this cache for second and later bounces, or as an indirect diffuse estimate,
while direct lighting remains sampled normally. This follows the NRC/GSCache
spirit and is likely the best first experiment.

\subsection{Formulation B: 3D Sample-Splatting Denoiser}

Instead of using Gaussians as scene primitives, treat each path sample or small
cluster of path samples as a temporary Gaussian in world space. Reproject/splat
those sample Gaussians into the current camera and nearby frames. This is the
3D analog of sample-based kernel splatting.

Advantages:

\begin{itemize}
  \item easier to prototype than a persistent learned cache,
  \item naturally uses path-sample data,
  \item can be compared directly to SVGF and KSN-like thinking,
  \item less likely to overfit long-term stale lighting.
\end{itemize}

Disadvantages:

\begin{itemize}
  \item may become just an expensive bilateral filter in disguise,
  \item memory bandwidth could dominate,
  \item without persistent learning it may not beat temporal accumulation.
\end{itemize}

\subsection{Formulation C: Gaussian Control Variate}

For a Monte Carlo integral
\[
  I = \int f(x)\,dx,
\]
if we can learn an approximation $g_\theta(x)$ whose integral is known or
cheaply estimated, then:
\[
  I =
  \int g_\theta(x)\,dx
  + \int (f(x)-g_\theta(x))\,dx.
\]
A corrected estimator is:
\[
  \hat{I} =
  G_\theta
  + \frac{1}{N}\sum_{i=1}^N
    \frac{f(x_i)-g_\theta(x_i)}{p(x_i)}.
\]

This is the cleanest route to unbiasedness. The issue is that a 3DGS radiance
cache will rarely have an analytic integral over the full path domain. A
practical compromise is to use the Gaussian cache as a biased high-order bounce
approximation and occasionally estimate/correct residuals. For vkGSplat's SDG
target, controlled bias may be acceptable if downstream task accuracy improves
and high-spp reference error is measured.

\subsection{Formulation D: Gaussian Path Guiding}

This is not denoising, but it may be more theoretically natural. Use Gaussians
to model spatial-directional incident radiance distributions, then sample
directions from the cache. This reduces variance before denoising. The path
guiding literature already uses mixture models and neural models; Gaussians are
a plausible explicit representation, especially if tied to primitive IDs and
surface tangent frames.

\section{Concrete vkGSplat Experiment Plan}

\subsection{Stage 0: Strengthen the Baseline}

Use the existing ray-tracing seed path plus current SVGF-style denoiser as the
baseline. Produce a fixed benchmark matrix:

\begin{itemize}
  \item scenes: Cornell box, Sponza/Bistro-like asset if available, synthetic
        robotics clutter scene,
  \item sample counts: 1, 2, 4, 8 spp,
  \item references: 512 or 1024 spp offline reference,
  \item metrics: MSE, relMSE, SSIM, FLIP if easy, temporal error, ghosting,
  \item costs: frame time, memory, update time, query time.
\end{itemize}

\subsection{Stage 1: Path Sample Logging}

Add an optional debug path-sample stream. Keep it CPU-first if necessary:

\begin{verbatim}
struct PathSampleRecord {
    float3 hit_position;
    float3 normal;
    float3 wo;
    float3 wi;
    float3 contribution;
    float3 throughput;
    float3 attenuated_radiance;
    float pdf;
    float roughness;
    uint32_t primitive_id;
    uint32_t path_depth;
    uint32_t sample_index;
    uint32_t pixel_x;
    uint32_t pixel_y;
};
\end{verbatim}

This is the critical enabling step. Without per-sample records, 3DGS can only
operate on pixels and loses most of its advantage.

For the first prototype, \texttt{attenuated\_radiance} and
\texttt{path\_depth} are the minimum GSCache-like fields. The other fields make
the cache safer for surfaces: normal, primitive ID, roughness, and direction
enable gating so radiance does not smear across unrelated transport events.

\subsection{Stage 2: CPU Gaussian Cache Prototype}

Implement a non-renderer intrusive prototype:

\begin{itemize}
  \item hash hit positions into a 3D grid,
  \item create/update one or more Gaussians per cell,
  \item align covariance to surface tangent/normal frame,
  \item store low-order spherical harmonics for directional radiance,
  \item store primitive ID histogram or dominant primitive ID,
  \item query cache for indirect radiance after first bounce,
  \item compare against path-traced reference and existing SVGF output.
\end{itemize}

The first version can be intentionally crude. The research question is whether
world-space cache reuse beats image-space history at equal cost.

A minimal host-side API could look like:

\begin{verbatim}
struct GaussianRadianceCacheOptions {
    float cell_size;
    float normal_reject_cosine;
    float roughness_bucket_width;
    uint32_t max_path_depth;
    uint32_t max_gaussians_per_cell;
};

struct GaussianRadianceQuery {
    float3 position;
    float3 normal;
    float3 wo;
    float roughness;
    uint32_t primitive_id;
    uint32_t path_depth;
};

class GaussianRadianceCache {
public:
    void reset(uint32_t width, uint32_t height);
    void update(span<const PathSampleRecord> samples);
    float3 query_indirect(const GaussianRadianceQuery& query) const;
    DenoiseFrame make_feature_frame(const ReprojectionFrame& current) const;
};
\end{verbatim}

The point of this API is to keep the experiment outside the core renderer until
there is evidence. It can consume path logs, produce feature frames, and plug
into the existing reprojection/denoise tests without committing vkGSplat to a
new backend architecture.

\subsection{Stage 3: Hybrid SVGF + Gaussian Cache}

Do not replace SVGF immediately. Feed the Gaussian cache output into the current
denoising pipeline as:

\begin{itemize}
  \item an auxiliary feature buffer,
  \item an indirect radiance prior,
  \item a variance estimate,
  \item a history validation signal.
\end{itemize}

This is likely to produce cleaner wins than a monolithic new denoiser.

\subsection{Stage 4: CUDA/Compute Path}

If Stage 2/3 shows signal, move to CUDA:

\begin{itemize}
  \item bin samples into tiles or hashed cells,
  \item update Gaussian moments with atomics or segmented reductions,
  \item compact/prune low-confidence Gaussians,
  \item query in a fixed neighborhood for predictable latency,
  \item expose cache output as a standard vkGSplat denoise input buffer.
\end{itemize}

\section{Possible Paper Angle}

A plausible paper title:

\begin{quote}
\emph{Gaussian Path-Space Caches for Compute-First Monte Carlo Reconstruction}
\end{quote}

The thesis would be:

\begin{quote}
For synthetic data generation, where scenes are controlled, camera/state
metadata is exact, and downstream model utility matters more than perfect
human-display latency, explicit Gaussian path-space caches provide a practical
middle ground between screen-space denoisers and neural radiance caches.
\end{quote}

This would fit vkGSplat well because the project is already about moving
rendering and reconstruction onto tensor-rich compute rather than defending
fixed-function graphics hardware.

\section{Evaluation Questions}

\begin{enumerate}
  \item Does a Gaussian cache reduce error at fixed samples per pixel compared
        with the current SVGF baseline?
  \item Does it reduce temporal instability under camera motion?
  \item Does it leak light across thin occluders or primitive boundaries?
  \item Does primitive-ID gating prevent most of the leaks?
  \item How stale does the cache become under moving lights or moving objects?
  \item Can low-order SH represent directional glossy transport well enough?
  \item Is the method still useful if only applied to diffuse/rough indirect
        lighting?
  \item Does it help downstream robotics perception metrics even when image
        metrics show small bias?
\end{enumerate}

\section{Current Recommendation}

Proceed, but with a narrow first target:

\begin{quote}
Build a Gaussian radiance cache for diffuse/rough indirect illumination in
static scenes, trained from path samples and evaluated against the existing
SVGF baseline. Treat it as a biased variance-reduction cache first. Only chase
unbiased control-variate theory after the practical signal is visible.
\end{quote}

The highest-value immediate engineering task is not a renderer rewrite. It is
path-sample logging plus a CPU cache prototype that can be run against existing
low-spp and high-spp test scenes.

\section{Overnight Research TODO}

\begin{itemize}
  \item Read the full GSCache paper after embargo/access allows, especially its
        cache hierarchy, update rules, and comparison to NRC.
  \item Check whether any 2026 arXiv papers cite GSCache or combine 3DGS with
        path guiding/control variates.
  \item Search for Gaussian mixture path guiding papers that can inform the
        directional component of the cache.
  \item Look for open-source implementations of StochasticSplats, 3DGS-MCMC,
        and DenoiseGS to borrow data structures or update heuristics.
  \item Sketch an interface for \path{GaussianRadianceCache} that does not
        disturb the current denoise API until evidence is stronger.
\end{itemize}

\begin{thebibliography}{99}

\bibitem{kerbl2023gs}
Bernhard Kerbl, Georgios Kopanas, Thomas Leimkuehler, and George Drettakis.
\newblock 3D Gaussian Splatting for Real-Time Radiance Field Rendering.
\newblock ACM Transactions on Graphics, 42(4), 2023.
\newblock \url{https://repo-sam.inria.fr/fungraph/3d-gaussian-splatting/}

\bibitem{kheradmand2024mcmc}
Shakiba Kheradmand, Daniel Rebain, Gopal Sharma, Weiwei Sun, Yang-Che Tseng,
Hossam Isack, Abhishek Kar, Andrea Tagliasacchi, and Kwang Moo Yi.
\newblock 3D Gaussian Splatting as Markov Chain Monte Carlo.
\newblock NeurIPS 2024; arXiv:2404.09591.
\newblock \url{https://arxiv.org/abs/2404.09591}

\bibitem{kheradmand2025stochasticsplats}
Shakiba Kheradmand, Delio Vicini, George Kopanas, Dmitry Lagun,
Kwang Moo Yi, Mark Matthews, and Andrea Tagliasacchi.
\newblock StochasticSplats: Stochastic Rasterization for Sorting-Free 3D
Gaussian Splatting.
\newblock arXiv:2503.24366, 2025.
\newblock \url{https://arxiv.org/abs/2503.24366}

\bibitem{zhou2024unified}
Yang Zhou, Songyin Wu, and Ling-Qi Yan.
\newblock Unified Gaussian Primitives for Scene Representation and Rendering.
\newblock arXiv:2406.09733, 2024.
\newblock \url{https://mangosister.github.io/gsr_site/}

\bibitem{bauer2026gscache}
David Bauer, Qi Wu, Hamid Gadirov, and Kwan Liu Ma.
\newblock GSCache: Real-Time Radiance Caching for Volume Path Tracing Using
3D Gaussian Splatting.
\newblock IEEE Transactions on Visualization and Computer Graphics,
32(1):901--911, 2026.
\newblock \url{https://research.rug.nl/en/publications/gscache-real-time-radiance-caching-for-volume-path-tracing-using-/}

\bibitem{schied2017svgf}
Christoph Schied, Anton Kaplanyan, Chris Wyman, Anjul Patney,
Chakravarty R. Alla Chaitanya, John Burgess, Shiqiu Liu, Carsten Dachsbacher,
Aaron Lefohn, and Marco Salvi.
\newblock Spatiotemporal Variance-Guided Filtering: Real-Time Reconstruction
for Path-Traced Global Illumination.
\newblock High Performance Graphics, 2017.
\newblock \url{https://research.nvidia.com/labs/rtr/publication/schied2017spatiotemporal/}

\bibitem{bako2017kpcn}
Steve Bako, Thijs Vogels, Brian McWilliams, Mark Meyer, Jan Novak,
Alex Harvill, Pradeep Sen, Tony DeRose, and Fabrice Rousselle.
\newblock Kernel-Predicting Convolutional Networks for Denoising Monte Carlo
Renderings.
\newblock ACM Transactions on Graphics, 36(4), 2017.
\newblock \url{https://civc.ucsb.edu/graphics/Papers/SIGGRAPH2017_KPCN/}

\bibitem{gharbi2019ksn}
Michael Gharbi, Tzu-Mao Li, Miika Aittala, Jaakko Lehtinen, and Fredo Durand.
\newblock Sample-Based Monte Carlo Denoising Using a Kernel-Splatting Network.
\newblock ACM Transactions on Graphics, 38(4), 2019.
\newblock \url{https://research.aalto.fi/en/publications/sample-based-monte-carlo-denoising-using-a-kernel-splatting-netwo}

\bibitem{muller2020ncv}
Thomas Mueller, Fabrice Rousselle, Alexander Keller, and Jan Novak.
\newblock Neural Control Variates.
\newblock ACM Transactions on Graphics, SIGGRAPH Asia, 2020.
\newblock \url{https://research.nvidia.com/labs/rtr/publication/muller2020ncv/}

\bibitem{muller2021nrc}
Thomas Mueller, Fabrice Rousselle, Jan Novak, and Alexander Keller.
\newblock Real-Time Neural Radiance Caching for Path Tracing.
\newblock ACM Transactions on Graphics, SIGGRAPH, 2021.
\newblock \url{https://research.nvidia.com/labs/rtr/publication/muller2021nrc/}

\bibitem{hasselgren2022mcd}
Jon Hasselgren, Nikolai Hofmann, and Jacob Munkberg.
\newblock Shape, Light, and Material Decomposition from Images Using Monte
Carlo Rendering and Denoising.
\newblock NeurIPS, 2022.
\newblock \url{https://research.nvidia.com/publication/2022-10_shape-light-and-material-decomposition-images-using-monte-carlo-rendering-and}

\bibitem{cheng2025denoisegs}
Yongsen Cheng, Yuanhao Cai, and Yulun Zhang.
\newblock DenoiseGS: Gaussian Reconstruction Model for Burst Denoising.
\newblock arXiv:2511.22939, 2025.
\newblock \url{https://arxiv.org/abs/2511.22939}

\bibitem{jiang2026denoisesplat}
Fuzhen Jiang, Zhuoran Li, and Yinlin Zhang.
\newblock DenoiseSplat: Feed-Forward Gaussian Splatting for Noisy 3D Scene
Reconstruction.
\newblock arXiv:2603.09291, 2026.
\newblock \url{https://arxiv.org/abs/2603.09291}

\bibitem{xu2026denoisegs}
Yabo Xu, Jin Ding, Jianbin Zhang, Ping Tan, and Mingrui Li.
\newblock Denoise-GS: Self-Supervised Denoising for Sparse-View 3D Gaussian
Splatting.
\newblock Sensors, 26(2):651, 2026.
\newblock \url{https://www.mdpi.com/1424-8220/26/2/651}

\end{thebibliography}

\end{document}
